\scanpage{76}
ramené au cas où c'est la clôture alg. de $k$, Mais alors \deleted{l'image de $S'$ dans $S$ est une partie \uline{fermée} et stable par générisation} le \inserted{saturé de $U'$ dans $S'$, étant une réunion d'ouverts,} \inserted{(les transformés de $U'$ par $\operatorname{Gal}(k'/k)$),} est un ouvert, et étant l'image inverse de l'image directe (fermée) \inserted{$U$} de $U'$ dans $S'$, est aussi fermée, ce qui implique que $U$ est \inserted{aussi} ouvert et fermé, et comme son image inverse sur $T$ est $\varnothing$, $U$ est $\varnothing$, donc $U'=\varnothing$.

cf simplification \unsure{éventuelle|éventuelles} de EGA \uline{IV} 4.

[\uline{Corollaire} 1 Supposons que $f:T\to S$ satisfasse $h_0$, et que $R'$ sur $k$ soit $\ne\varnothing$. Alors le diagramme
\diagram{\begin{tikzcd}[ampersand replacement=\&]
\mathrm{OF}(S)\arrow[r,hook]\arrow[d] \& \mathrm{OF}(T)\arrow[d]\\
\mathrm{OF}(S')\arrow[r,hook] \& \mathrm{OF}(T')
\end{tikzcd}}
est \uline{cartésien}.

En effet, $\mathrm{OF}(S)\unsure{-|\to}\mathrm{OF}(T)$ est injectif, il en est de $\overline m$ à fortiori de $\mathrm{OF}(S)\to\mathrm{OF}(T)\times_{\mathrm{OF}(T')}\mathrm{OF}(S')$, prouvons que c'est \uline{surjectif}. Soit donc $V\in\mathrm{OF}(T)$ dont l'image dans $\mathrm{OF}(T')$ provient de $U'\in\mathrm{OF}(S')$, prouvons que $U'$ provient de $U\in\mathrm{OF}(S)$

\uline{I}mmédiat par descente, en utilisant que $\mathrm{OF}(S'')\unsure{-|\to}\mathrm{OF}(T'')$ est injectif\ldots\ .
% archive sheet number: 75
